% Part:sets-functions-relations
% Chapter: size-of-sets
% Section: pairing-alt

\documentclass[../../../include/open-logic-section]{subfiles}

\begin{document}

\olfileid{sfr}{siz}{pai-alt}

\olsection{Een andere paarfunctie}

\begin{explain}
Er zijn andere opsommingen van $\Nat^2$ waarvan de inverse gemakkelijker te
bepalen is. Hier volgt er een. Geef de opsomming niet weer in een matrix, maar
begin met de lijst van positieve gehele getallen en zet bij elk getal een
aanvankelijk leeg vak. Stel dat we deze vakken achtereenvolgens als volgt met
paren $\tuple{n,m}$ vullen. Begin met de paren die op de eerste plaats~$0$
hebben (dus de paren $\tuple{0,m}$). Zet het eerste paar (dus
$\tuple{0,0}$) in het eerste lege vak, sla daarna een leeg vak over, zet het
tweede paar (dus $\tuple{0,1}$) in het volgende lege vak, sla opnieuw een vak
over, enzovoort. Het (onvolledige) begin van onze opsomming ziet er nu zo uit:
\[\small
\begin{array}{@{}c c c c c c c c c c c@{}}
\mathbf 1 & \mathbf 2 & \mathbf 3 & \mathbf 4 & \mathbf 5 & \mathbf 6 & \mathbf 7 & \mathbf 8 & \mathbf 9 & \mathbf{10} & \dots \\ \\
\tuple{0,0} &  & \tuple{0,1} &  & \tuple{0,2} &  & \tuple{0,3} & & \tuple{0,4} &  & \dots \\
\end{array}
\]
Herhaal dit voor de nog lege vakken met de paren $\tuple{1,m}$ en sla daarbij
opnieuw steeds één leeg vak over:
\[\small
\begin{array}{@{}c c c c c c c c c c c@{}}
\mathbf 1 & \mathbf 2 & \mathbf 3 & \mathbf 4 & \mathbf 5 & \mathbf 6 & \mathbf 7 & \mathbf 8 & \mathbf 9 & \mathbf{10} & \dots \\ \\
\tuple{0,0} & \tuple{1,0} & \tuple{0,1} &  & \tuple{0,2} & \tuple{1,1} & 
\tuple{0,3} & & \tuple{0,4} &  \tuple{1,2} & \dots \\
\end{array}
\]
Vul op dezelfde manier de paren $\tuple{2,m}$, $\tuple{3,m}$, enzovoort in.
Onze voltooide opsomming begint dan als volgt:
\[\small
\begin{array}{@{}cc c c c c c c c c c@{}}
\mathbf 1 & \mathbf 2 & \mathbf 3 & \mathbf 4 & \mathbf 5 & \mathbf 6 & \mathbf 7 & \mathbf 8 & \mathbf 9 & \mathbf{10} & \dots \\ \\
\tuple{0,0} & \tuple{1,0} & \tuple{0,1} & \tuple{2,0}  & \tuple{0,2} & 
\tuple{1,1} & \tuple{0,3} & \tuple{3,0}  & \tuple{0,4} &  \tuple{1,2} & \dots \\
\end{array}
\]
Als we de vakken van de bovenstaande matrix volgens deze opsomming nummeren,
krijgen we geen keurige zigzaglijn, maar de volgende rangschikking:
\[
\begin{array}{ c | c | c | c | c | c | c | c }
& \mathbf 0 & \mathbf 1 & \mathbf 2 & \mathbf 3 & \mathbf 4 & \mathbf 5 & \dots \\
\hline
\mathbf 0 & 1 & 3 & 5 & 7 & 9 & 11 & \dots \\
\hline
\mathbf 1 & 2 & 6 & 10 & 14 & 18 & \dots & \dots \\
\hline
\mathbf 2 & 4 & 12 & 20 & 28 & \dots & \dots & \dots \\
\hline
\mathbf 3 & 8 & 24 & 40 & \dots & \dots & \dots & \dots \\
\hline
\mathbf 4 & 16 & 48 & \dots & \dots & \dots & \dots & \dots \\
\hline
\mathbf 5 & 32 & \dots & \dots & \dots & \dots & \dots & \dots \\
\hline
\vdots & \vdots & \vdots & \vdots & \vdots & \vdots & \vdots & \ddots\\
\end{array}
\]
We zien dat de paren in rij~$0$ op de oneven genummerde plaatsen van onze
opsomming staan: het paar $\tuple{0,m}$ staat op plaats $2m+1$. De paren in
de tweede rij, $\tuple{1,m}$, staan op plaatsen waarvan het nummer het dubbele
van een oneven getal is, om precies te zijn $2 \cdot (2m+1)$. De paren in de
derde rij, $\tuple{2,m}$, staan op plaatsen waarvan het nummer viermaal een
oneven getal is, $4 \cdot (2m+1)$; enzovoort. De factoren van $(2m+1)$ voor
de opeenvolgende rijen, $1$, $2$, $4$, $8$, \dots, zijn precies de machten
van~$2$: $1= 2^0$, $2 = 2^1$, $4 = 2^2$, $8 = 2^3$, \dots\@ De bijbehorende
exponent is steeds het eerste lid van het betreffende paar. Voor het paar
$\tuple{n,m}$ is de factor dus $2^n$. Dit levert de algemene formule
$2^n \cdot (2m+1)$ op. Deze afbeelding stuurt paren echter naar
\emph{positieve} gehele getallen: $\tuple{0,0}$ staat op plaats~$1$. Willen we
bij plaats~$0$ beginnen, dan moeten we van het resultaat~$1$ aftrekken. Dit
geeft:
\end{explain}

\begin{ex}
De functie $h\colon \Nat^2 \to \Nat$ die wordt gegeven door
\[
h(n,m) = 2^n (2m+1) - 1
\]
is een paarfunctie voor de verzameling paren natuurlijke getallen~$\Nat^2$.
\end{ex}

\begin{explain}
In onze tweede opsomming van $\Nat^2$ heeft het paar $\tuple{0,0}$ dus de code
$h(0,0) = 2^0(2\cdot 0+1) - 1 = 0$; het paar $\tuple{1,2}$ heeft de code
$2^{1} \cdot (2 \cdot 2 + 1) - 1 = 2 \cdot 5 - 1 = 9$; en het paar
$\tuple{2,6}$ heeft de code $2^{2} \cdot (2 \cdot 6 + 1) - 1 = 51$.
\end{explain}

Soms volstaat het de paren natuurlijke getallen~$\Nat^2$ te coderen zonder te
eisen dat de codering surjectief is. De inversen van zulke coderingen zijn
slechts partiële functies.

\begin{ex}
De functie $j\colon \Nat^2 \to \Nat^+$ die wordt gegeven door
\[
j(n,m) = 2^n3^m
\]
is !!a{injective} functie $\Nat^2 \to \Nat$.
\end{ex}
\end{document}